% !PRINTED 36
\weberpage{36}

The first step is to form the equation
\begin{equation}
 A={a_0\over b_0}B+C,
\tag{6}
\end{equation}
where
\begin{equation}
 C=c_0x+c_1
\tag{7}
\end{equation}
and
\begin{equation}
 c_0={a_1b_0-b_1a_0\over b_0},
 \qquad
 c_1={a_2b_0-b_2a_0\over b_0}.
\tag{8}
\end{equation}
If $c_0=0$, so that $C$ is constant, then $A,B$ have a common factor only if also $c_1=0$; and then $A$ is divisible by $B$, that is, $A$ and $B$ differ only by a constant factor. If however $c_0$ is different from zero, we continue the calculation by putting
\[
 B=QC+D
\]
where, according to \S~4, if there $\alpha=-c_1:c_0$ is put,
\begin{equation}
 D={b_0c_1^2-b_1c_0c_1+b_2c_0^2\over c_0^2}.
\tag{9}
\end{equation}
If $D$ is non-zero, then $A,B$ have no common divisor. If however $D=0$, then $C$ is the greatest common factor of $A$ and $B$.

Substituting the values $c_0,c_1$ from (8) into (9), and suppressing the denominator $b_0c_0^2$, one obtains the condition for the existence of a common divisor $C$ of $A$ and $B$ in the form
\begin{equation}
\begin{aligned}
&a_0^2b_2^2+a_2^2b_0^2-2a_0a_2b_0b_2-a_1a_2b_0b_1-a_0a_1b_1b_2\\
&\qquad{}+a_0a_2b_1^2+a_1^2b_0b_2=0
\end{aligned}
\tag{10}
\end{equation}
or
\begin{equation}
 (a_0b_2-b_0a_2)^2+(a_0b_1-a_1b_0)(a_2b_1-a_1b_2)=0.
\tag{11}
\end{equation}
This condition is also fulfilled when $c_0$ and $c_1$ are both zero, and is therefore the necessary and sufficient condition that $A,B$ have a common divisor.

The left hand side of (10) or (11) is called the \gesperrt{resultant of the functions $A$ and $B$}.

As a second example we take the one already chosen in \S~3. Put
\begin{equation}
\begin{aligned}
 f(x)&=a_0x^3+a_1x^2+a_2x+a_3=A,\\
 f'(x)&=3a_0x^2+2a_1x+a_2=B,
\end{aligned}
\tag{12}
\end{equation}
and suppose $a_0$ different from zero. Then by \S~3 we have:
%CONT>>
